\label{sec:borders}
\begin{proposition}[Two Bratteli diagrams of the depth tower]\label{prop:twodiagrams}
\begin{enumerate}[leftmargin=*]
\item \emph{Tile diagram.} The vertices at level $l$ are the nonempty depth-$l$ tiles, edges are inclusions, and $\nu(v)=\gamma(v)$ is a one-sided
state ($\nu(v)=\sum_{\text{children}}\nu$). It is a tree, so its AF algebra is commutative. On a great circle, with the angular order as
$\le_s$, it is Putnam--Trevi\~no's ``decimal expansion'' picture: the circle is the quotient of the path space by successor
identifications, and the successor pairs are the two sides of the breakpoints.
\item \emph{Fiber diagram.} The supertiles of level $l$ are the fibres of $h_l$; write $x\approx_lx'$ iff $h_l(x)=h_l(x')$. These relations
increase with $l$ (the AF direction), their classes are polyhedra of dimension $n-\operatorname{rank}(D_lW_l\cdots D_1W_1)$ (about $n/2$ from the
first layer on), and $\approx_\infty=\bigcup_l\approx_l$ (eventual indistinguishability) is contained in the tail relation of the \emph{letter
diagram}, whose level-$l$ vertices are the layer-$l$ gate patterns alone (finite rank $2^n$). The letter diagram's tail relation is the network's
$R_{AF}$.
\item \emph{Borders.} Call $n-j$ the border dimension of $x$ if, for every large depth, $x$ lies on a face of codimension $j$. At infinite width
and infinite depth every Gaussian ball, of any radius $\sqrt\hbar>0$, meets walls of every layer, with $\sum_ln\,\theta_{l-1}(\hbar)/2\pi=\infty$ by
Theorem~\ref{thm:fatou}. So points of border dimension $<n$ are dense; they are null. This is Julien--Savinien's thin set (their
Prop.~3.22 and Rem.~3.26), and it is dense precisely because depth is parabolic: under the hyperbolic finite-width dynamics of
Conjecture~\ref{conj:hyper} the sum converges.
\item \emph{Births.} A wall of layer $l$ inside a tile of depth $l-1$ is a face that does not descend from a face of the parent. The face
multi-diagram (Bratteli diagrams of faces of each codimension, joined by incidence edges) therefore has \emph{birth} edges from tiles at
level $l-1$ to codimension-one faces at level $l$. These are the reverse of Julien--Savinien's escaping edges (faces swallowed by
supertiles), and they are where stage 10 locates the creation of contacts.
\end{enumerate}
\end{proposition}
\begin{proof}
(1) and (2) are immediate from the definitions; $h_l(x)=h_l(x')$ forces equal gates at all later layers. (3) uses the infinite-width
unresolved variance at every layer and the divergence of $\sum\theta_l$. (4) is the definition of a wall.
\end{proof}

\begin{proposition}[The mean is a border functional]\label{prop:border}
For any network output $f$ (piecewise linear, $1$-homogeneous),
\[
\E_\gamma f=\E\Delta f=\frac1{\sqrt{2\pi}}\sum_{F}\kappa_F\,\alpha_F ,
\]
where $F$ runs over the codimension-one faces of the depth-$L$ tiling, $\kappa_F$ is the jump of the normal derivative across $F$, and $\alpha_F$
is the Gaussian measure of the cone $F$ inside its hyperplane (its solid-angle fraction). Equivalently, let $H_\partial=L^2(\bigcup_FF\times\{+,-\},\varphi\,dS)$
(two copies of each wall piece, one per side), let $\Gamma$ be the side grading, $\mathbb F$ the side flip, and $\pi(g)$ multiplication by the
one-sided values of a field $g$ that is constant on open tiles. Then
\[
[\mathbb F,\pi(g)]=0\iff g\ \text{is continuous across every wall},\qquad
\E_\gamma f=\tfrac12\Tr_\varphi\big(\Gamma\,\mathbb F\,[\mathbb F,\pi(\partial_\nu f)]\big).
\]
\end{proposition}
\begin{proof}
Euler and Stein give $\E f=\E[x\cdot\nabla f]=\E\Delta f$, and $\Delta f=\sum_F\kappa_F\,\delta_F$ as a distribution. Moreover $\int_F\varphi_n\,dS=\varphi_1(0)\,\alpha_F$
because $F$ is a cone through $0$. On one face, with $\pi(g)=\diag(g_+,g_-)$, $\Gamma\mathbb F[\mathbb F,\pi(g)]=\diag(g_+-g_-,g_+-g_-)$; take $g=\partial_\nu f$.
\end{proof}
The second form is Putnam--Trevi\~no's Fredholm module with the successor flip replaced by the wall flip. Their Theorem 9.7 (functions
descend to the quotient iff they commute with the flip) becomes continuity across walls, and the Gaussian mean of the network becomes the
supertrace pairing of the flip with the normal derivative: the kink current of stage 10. On a continuum of wall points the commutators
are multiplication operators rather than compact ones, so this is a Gaussian-weighted analogue, not a Fredholm module in the strict
sense. Its restriction to a great circle, which has finitely many breakpoints, is one. The two-dimensional check (Verification,
item 7) has $70$ breakpoints and gives $\E f=0.36919251$ both by angular quadrature and by the border sum. In Julien--Savinien's terms, the
whole mean is carried by the border paths of the codimension-one face diagram: the generic (AF) points contribute nothing.
